\begin{PersianAbstract}
	
	``
	دیستروفی عضلانی دوشن 
	\footnote{\lr{Duchenne Muscular Dystrophy (DMD)}} 
	یک بیماری عصبی‌ـ‌عضلانی پیشرونده است که با ضعف تدریجی عضلات اسکلتی و کاهش توانایی‌های حرکتی همراه می‌شود. مقیاس ارزیابی سرپایی نورث‌استار (
	 \lr{NSAA}\footnote{\lr{North Star Ambulatory Assessment}}
	) یکی از ابزارهای متداول برای ارزیابی عملکرد حرکتی بیماران سرپایی مبتلا به این بیماری است و برای هر فعالیت، بر اساس مشاهده بالینی، نمره‌ای از صفر تا دو اختصاص می‌دهد. وابستگی این فرایند به قضاوت ارزیاب، ضرورت بررسی روش‌های کمی و تکرارپذیر برای تحلیل نحوه اجرای حرکات را مطرح می‌کند. هدف پژوهش حاضر، بررسی امکان‌پذیری فنی طبقه‌بندی سطوح عملکردی شبیه‌سازی‌شده متناظر با نمرات فعالیت‌های منتخب 
	\lr{NSAA}
	 با استفاده از حسگرهای اینرسیایی پوشیدنی و الگوریتم‌های یادگیری ماشین بود.
	در این پژوهش، داده‌برداری مستقیم از بیماران مبتلا به دیستروفی عضلانی دوشن انجام نشد. داده‌های مورد استفاده توسط یک فرد سالم آشنا با بیماری و معیارهای 
	\lr{NSAA}
	 تولید شدند که سه سطح عملکردی را در شرایط کنترل‌شده شبیه‌سازی کرد. پنج فعالیت شامل ایستادن، برخاستن از صندلی، پرش با دو پا، بالا آوردن سر و بالا رفتن از جعبه با پای راست بررسی شدند. پنج حسگر اینرسیایی روی پیشانی، مچ‌های دست و مچ‌های پا نصب شدند و شتاب خطی و سرعت زاویه‌ای را با نرخ تقریبی ۹۰ هرتز ثبت کردند. پس از اعمال فیلتر پایین‌گذر باترورث مرتبه چهارم با فرکانس قطع ۴ هرتز و آماده‌سازی سیگنال‌ها، از هر کانال ۲۳ ویژگی حوزه زمان استخراج شد که در مجموع برداری شامل ۶۹۰ ویژگی برای هر نمونه ایجاد کرد. استانداردسازی و تحلیل مؤلفه‌های اصلی درون یک خط پردازش انجام شدند. سپس عملکرد ماشین بردار پشتیبان، نزدیک‌ترین همسایه و جنگل تصادفی با استفاده از اعتبارسنجی متقاطع طبقه‌بندی‌شده پنج‌تایی و تکرارشونده مقایسه شد.
	نتایج نشان داد که ماشین بردار پشتیبان با هسته شعاعی و تحلیل مؤلفه‌های اصلی برای هر پنج فعالیت مناسب‌ترین عملکرد را ارائه کرد. بسته به نوع فعالیت، ۱۰ تا ۱۵ مؤلفه اصلی برای دستیابی به عملکرد پایدار انتخاب شدند. مدل‌های منتخب به حد نصاب داخلی دقت ۹۰ درصد برای فعالیت‌های ایستادن، پریدن و بالا آوردن سر و حد نصاب ۸۰ درصد برای برخاستن از صندلی و بالا رفتن از جعبه دست یافتند. این مقادیر معیارهای مهندسی داخلی برای انتخاب مدل هستند و نباید به‌عنوان شاخص اعتبار بالینی سامانه تفسیر شوند.
	یافته‌ها امکان‌پذیری فنی تفکیک الگوهای حرکتی شبیه‌سازی‌شده را در داده‌های یک فرد نشان می‌دهند، اما درباره عملکرد سامانه در بیماران واقعی یا تعمیم آن به افراد جدید نتیجه‌گیری نمی‌کنند. ارزیابی بالینی سامانه مستلزم جمع‌آوری داده از بیماران، امتیازدهی هم‌زمان توسط متخصصان و اعتبارسنجی مستقل از آزمودنی است.
	``
	
	\keywords{
		دیستروفی عضلانی دوشن، حسگر اینرسیایی، تحلیل حرکت، یادگیری ماشین، مقیاس 
		\lr{NSAA}
		، ارزیابی عملکرد حرکتی
	}
	
	
\end{PersianAbstract}
